\section*{Problemas}

\subsection*{Enunciados de los problemas}

% Recordar
\begin{problema}
	\label{prob:71ce5a0}
	Enuncie, sin realizar ningún cálculo, la función de masa de probabilidad de un ensayo de Bernoulli con parámetro $p$, y escriba de memoria las fórmulas de $\mathbb{E}[X]$ y $\text{Var}(X)$ para $X \sim \text{Bin}(n,p)$.
\end{problema}

% Comprender
\begin{problema}
	\label{prob:1149be6}
	Explique, sin invocar la fórmula binomial de memoria, por qué el número de secuencias de $n$ ensayos Bernoulli independientes que producen exactamente $k$ éxitos es $\comb{n}{k}$, y utilice ese argumento para justificar por qué $P(Y=k) = \comb{n}{k}p^k q^{n-k}$.
\end{problema}

% Aplicar
\begin{problema}
	\label{prob:c3d2032}
	Una firma de ciberseguridad monitorea $n=12$ intentos de acceso independientes, cada uno clasificado como malicioso con probabilidad $p=0.15$. Sea $X \sim \text{Bin}(12, 0.15)$ el número de accesos maliciosos detectados.
	\begin{enumerate}
		\item Calcule $P(X \le 2)$.
		\item Calcule la probabilidad condicional $P(X \ge 4 \mid X \ge 1)$.
	\end{enumerate}
\end{problema}

% Analizar
\begin{problema}
	\label{prob:bae56b2}
	Sean $X \sim \text{Bin}(n_1,p)$ y $Y \sim \text{Bin}(n_2,p)$ independientes, con el mismo parámetro de éxito $p$. Demuestre, usando convolución discreta y la Identidad de Vandermonde $\sum_{j=0}^m \comb{n_1}{j}\comb{n_2}{m-j}=\comb{n_1+n_2}{m}$, que $Z=X+Y \sim \text{Bin}(n_1+n_2,p)$, y verifique que la aditividad de $\mathbb{E}[\cdot]$ y $\text{Var}(\cdot)$ se conserva bajo esta suma.
\end{problema}

% Evaluar
\begin{problema}
	\label{prob:19f50da}
	Un ingeniero de control de calidad, ante $X \sim \text{Bin}(n=500, p=0.04)$, decide aproximar $X$ por una distribución Normal $N(np, np(1-p))$ para calcular la probabilidad de aceptar un lote (criterio $X \le 25$), sin verificar ninguna condición previa y sin aplicar corrección de continuidad. Evalúe esta decisión: ¿es válida la aproximación Normal en este caso? Justifique citando la regla operativa usual ($np \ge 5$ y $n(1-p) \ge 5$), y explique qué error introduce omitir la corrección de continuidad de Yates.
\end{problema}

% Crear
\begin{problema}
	\label{prob:3c14103}
	Diseñe un ejemplo original (contexto, valores de $n$ y $p$, y la variable aleatoria de interés) que modele un conteo de éxitos mediante una distribución Binomial, distinto de los ejemplos de control de calidad o ciberseguridad ya vistos en esta sección. Plantee explícitamente la pregunta de probabilidad que resolvería su ejemplo y calcule $\mathbb{E}[X]$ y $\text{Var}(X)$ para los parámetros que usted eligió.
\end{problema}

\subsection*{Soluciones detalladas}

% Recordar
\begin{solproblema}[prob:71ce5a0]
	La función de masa de probabilidad de un ensayo Bernoulli($p$) es $f_X(x) = p^x(1-p)^{1-x}$ para $x \in \set{0,1}$. Para $X \sim \text{Bin}(n,p)$, suma de $n$ ensayos Bernoulli independientes, $\mathbb{E}[X] = np$ y $\text{Var}(X) = np(1-p)$.
\end{solproblema}

% Comprender
\begin{solproblema}[prob:1149be6]
	Cada secuencia particular de $n$ ensayos con exactamente $k$ éxitos y $n-k$ fracasos tiene, por independencia, probabilidad $p^k q^{n-k}$ (el orden de los ensayos no afecta el producto). El número de formas distintas de elegir en cuáles de las $n$ posiciones ocurren los $k$ éxitos es el número de subconjuntos de tamaño $k$ de un conjunto de $n$ elementos, es decir $\comb{n}{k}$. Como estas secuencias son eventos mutuamente excluyentes que juntas forman el evento $\set{Y=k}$, sumar sus probabilidades individuales iguales da $P(Y=k) = \comb{n}{k}p^k q^{n-k}$.
\end{solproblema}

% Aplicar
\begin{solproblema}[prob:c3d2032]
	\begin{enumerate}
		\item Con $P(X=0)\approx0.14224$, $P(X=1)\approx0.30121$, $P(X=2)\approx0.29236$:
		\begin{align*}
			P(X\le2) = 0.14224+0.30121+0.29236 \approx 0.7358.
		\end{align*}
		\item Con $P(X=3)\approx0.17197$, se tiene $P(X\le3)\approx0.90778$. Entonces:
		\begin{align*}
			P(X\ge4\mid X\ge1) = \dfrac{1-P(X\le3)}{1-P(X=0)} = \dfrac{1-0.90778}{1-0.14224} = \dfrac{0.09222}{0.85776}\approx0.1075.
		\end{align*}
	\end{enumerate}
\end{solproblema}

% Analizar
\begin{solproblema}[prob:bae56b2]
	Por independencia, la probabilidad de que $Z=X+Y=m$ se obtiene por convolución discreta:
	\begin{align*}
		P(Z=m) = \sum_{j=0}^m P(X=j)P(Y=m-j) = \sum_{j=0}^m \comb{n_1}{j}p^j q^{n_1-j} \comb{n_2}{m-j}p^{m-j}q^{n_2-(m-j)} = p^m q^{(n_1+n_2)-m}\sum_{j=0}^m \comb{n_1}{j}\comb{n_2}{m-j}.
	\end{align*}
	Por la Identidad de Vandermonde, la suma remanente es $\comb{n_1+n_2}{m}$, de modo que $P(Z=m) = \comb{n_1+n_2}{m}p^m q^{(n_1+n_2)-m}$, es decir $Z \sim \text{Bin}(n_1+n_2,p)$. La aditividad se verifica directamente: $\mathbb{E}[Z]=(n_1+n_2)p=\mathbb{E}[X]+\mathbb{E}[Y]$ y $\text{Var}(Z)=(n_1+n_2)p(1-p)=\text{Var}(X)+\text{Var}(Y)$.
\end{solproblema}

% Evaluar
\begin{solproblema}[prob:19f50da]
	Con $np = 500(0.04) = 20 \ge 5$ y $n(1-p) = 500(0.96) = 480 \ge 5$, la regla operativa usual sí se satisface, por lo que la aproximación Normal es razonable en principio. Sin embargo, omitir la corrección de continuidad de Yates introduce un sesgo sistemático evitable: al aproximar una variable discreta por una continua, $P(X \le 25)$ debe calcularse como $\Phi\!\left(\dfrac{25.5-20}{\sqrt{19.2}}\right)$ y no como $\Phi\!\left(\dfrac{25-20}{\sqrt{19.2}}\right)$; ignorar el término $0.5$ subestima la probabilidad acumulada real en aproximadamente el área de la curva Normal entre $25$ y $25.5$. La decisión del ingeniero es parcialmente válida (las condiciones de aplicabilidad se cumplen) pero técnicamente incompleta por omitir la corrección de continuidad.
\end{solproblema}

% Crear
\begin{solproblema}[prob:3c14103]
	\textbf{Ejemplo:} un estudio agrícola siembra $n=15$ semillas de una variedad con probabilidad de germinación $p=0.7$ cada una, de forma independiente. Sea $X$ el número de semillas que germinan, $X \sim \text{Bin}(15, 0.7)$. Pregunta de interés: ¿cuál es la probabilidad de que germinen al menos $12$ semillas, $P(X \ge 12)$? Los momentos son $\mathbb{E}[X] = np = 15(0.7) = 10.5$ y $\text{Var}(X) = np(1-p) = 15(0.7)(0.3) = 3.15$. La probabilidad $P(X\ge12)$ se calcularía sumando $P(X=12)+P(X=13)+P(X=14)+P(X=15)$, cada término evaluado con $\comb{15}{k}(0.7)^k(0.3)^{15-k}$.
\end{solproblema}
